\chapter{RESULTS AND DISCUSSIONS}
\label{ChapterResults}

This chapter presents the Phase~1 implementation progress, preliminary outcomes, and an academic discussion of the proposed methodology. Results are reported conservatively, clearly distinguishing implemented components from those deferred to Phase~2.

---

\section{Major Findings from Literature}
\label{sec:major_findings}

The systematic literature review conducted in Chapter~2 yields the following key findings relevant to TemporalAML:

\begin{enumerate}
    \item \textbf{GCN-based static AML is insufficient for temporal laundering detection.} Weber et al.~\cite{weber2019anti} established GCN as a viable AML tool on the Elliptic dataset, but the static graph representation collapses transaction timestamps, enabling future data leakage and destroying velocity signals essential for detecting layering chains.

    \item \textbf{Discrete snapshot methods lose sub-interval temporal dynamics.} EvolveGCN~\cite{pareja2020evolvegcn} evolves model weights across graph snapshots but cannot capture continuous temporal velocity within a snapshot window—a critical limitation given that rapid laundering chains may complete within a single time step.

    \item \textbf{Fixed-frequency temporal encodings are insufficient for cryptocurrency.} TGAT~\cite{xu2020tgat} introduced continuous-time graph attention but uses geometrically fixed sinusoidal frequencies designed for NLP token distributions. Cryptocurrency transaction intervals follow bursty, non-stationary, power-law distributions that fixed frequencies cannot represent optimally.

    \item \textbf{Binary AML labels are inadequate for compliance operations.} All reviewed GNN-based AML methods produce binary illicit/licit scores. Regulatory compliance (FATF, FinCEN) requires the identification of specific laundering typologies and traceable evidence chains for SAR submission—a gap unaddressed by any single reviewed work.

    \item \textbf{Explainability is a necessary, not optional, AML capability.} GNNExplainer~\cite{ying2019gnnexplainer} and related work demonstrate the technical feasibility of subgraph evidence extraction, but no prior work has applied per-pattern, time-ordered explainability specifically to cryptocurrency AML detection.
\end{enumerate}

Table~\ref{tab:lit_summary_findings} summarises the method-gap correspondence from the literature review.

\begin{table}[!htbp]
\centering
\caption[Literature Finding Summary]{Key findings from literature and their correspondence to the TemporalAML research objectives.}
\label{tab:lit_summary_findings}
\begin{tabularx}{\textwidth}{>{\hsize=0.85\hsize}L >{\hsize=1.0\hsize}L >{\hsize=1.15\hsize}L}
\toprule
\textbf{Literature Finding} & \textbf{Identified Gap} & \textbf{TemporalAML Response} \\
\midrule
GCN/GAT: static graphs, temporal leakage & No chronological ordering or causal filtering & Causal temporal masking + chronological data splits (O1) \\
\addlinespace
TGAT: fixed sinusoidal frequencies & Cannot adapt to bursty crypto intervals & Learnable Fourier frequency parameters $\boldsymbol{\omega}$ (O2) \\
\addlinespace
EvolveGCN: discrete snapshots & Sub-interval velocity is lost & Continuous-time $\Delta t$ encoding in attention (O2) \\
\addlinespace
All GNN AML: binary labels only & No pattern identification; non-actionable & Three pattern-specific classification heads (O3) \\
\addlinespace
GNNExplainer: no per-pattern time-order & Binary or class-level explanations only & Per-pattern, time-ordered evidence subgraph (O4) \\
\bottomrule
\end{tabularx}
\end{table}

---

\section{Summary of the Proposed Design}
\label{sec:proposed_design_summary}

Table~\ref{tab:phase1_progress} presents the Phase~1 implementation progress summary for all four objectives.

\begin{table}[!htbp]
\centering
\caption[Phase 1 Implementation Progress]{TemporalAML Dissertation Phase 1 implementation progress across all research objectives.}
\label{tab:phase1_progress}
\small
\begin{tabularx}{\textwidth}{c >{\hsize=0.95\hsize}L >{\hsize=1.05\hsize}L c >{\hsize=1.0\hsize}L}
\toprule
\textbf{Obj.} & \textbf{Planned Work} & \textbf{Phase 1 Deliverable} & \textbf{Status} & \textbf{Key Artefact} \\
\midrule
O1 & Construct temporal graph from Elliptic dataset; engineer 6 topological features; leakage-free chronological split. & 172-dim feature matrix; PyG Data object; train/val/test split masks; zero-leakage unit tests. & \textbf{Completed} & \texttt{data\_loader.py}; \texttt{elliptic\_}\allowbreak\texttt{graph.pt}; \texttt{test\_no\_}\allowbreak\texttt{leakage.py} \\
\addlinespace
O2 & Design Learnable Fourier Time Encoder; integrate into TGAT attention; verify gradient flow. & \texttt{FourierTime}\allowbreak\texttt{Encoder}; \texttt{TGATConv} layer; causal temporal neighbour sampling. & \textbf{Completed} & \texttt{src/models.py} \\
\addlinespace
O3 & Mine pattern labels (Tarjan SCC, BFS, degree thresholding); build shared TGAT backbone with 3 pattern heads; joint multi-task training. & Pattern label mining; \texttt{MultiPattern}\allowbreak\texttt{TGAT} model; joint loss; preliminary training runs. & \textbf{In Progress} & \texttt{src/train.py}; \texttt{temporalaml\_}\allowbreak\texttt{best.pth} \\
\addlinespace
O4 & Integrate GNNExplainer; generate per-pattern, time-ordered subgraph evidence; design SAR narrative template. & \texttt{src/explain.py}; GNNExplainer integration; Streamlit forensic dashboard with subgraph visualisation. & \textbf{Implemented} & \texttt{src/explain.py}; \texttt{dashboard.py} \\
\bottomrule
\end{tabularx}
\end{table}

Table~\ref{tab:phase_scope} compares Phase~1 and Phase~2 scope.

\begin{table}[!htbp]
\centering
\caption[Phase 1 vs Phase 2 Scope]{Comparison of Phase 1 and Phase 2 dissertation scope.}
\label{tab:phase_scope}
\begin{tabularx}{\textwidth}{>{\hsize=1.1\hsize}L >{\hsize=0.9\hsize}L >{\hsize=1.0\hsize}L}
\toprule
\textbf{Activity} & \textbf{Phase 1 Status} & \textbf{Phase 2 Plan} \\
\midrule
Temporal Graph Construction (O1) & Completed & Not applicable \\
Learnable Fourier Encoding (O2) & Completed & Ablation study \\
Multi-Pattern TGAT (O3) & In progress (architecture implemented; training ongoing) & Full training; hyperparameter tuning \\
GNNExplainer Integration (O4) & Implemented & Quality evaluation (fidelity, sparsity) \\
Benchmarking (GCN, GraphSAGE, GNN-LSTM) & Not started & Comprehensive evaluation \\
F1, AUC-ROC, Precision, Recall metrics & Preliminary/partial & Full evaluation on test set \\
Ablation Study & Not started & Isolation of Fourier encoding contribution \\
SAR generation quality evaluation & Not started & Expert review \\
\bottomrule
\end{tabularx}
\end{table}

---

\section{Phase 1 Implementation Results}
\label{sec:phase1_results}

\subsection{Dataset Preparation and Graph Construction (O1)}

The following confirmed outcomes are reported for Objective~1:
\begin{itemize}
    \item The three Elliptic CSV files were successfully loaded and joined. All 203,769 transaction node IDs were remapped to contiguous indices $[0, 203768]$ via a bidirectional hash map.
    \item A directed edge index tensor $\mathbf{E} \in \mathbb{Z}^{2 \times 234355}$ was constructed.
    \item Six topological features were engineered and appended, yielding a final feature matrix $\mathbf{X} \in \mathbb{R}^{203769 \times 172}$.
    \item A \texttt{StandardScaler} was fitted exclusively on training-split nodes ($t \leq 34$), with transformation applied to all partitions.
    \item Chronological split masks were assigned: Train ($t \leq 34$: 143,551 nodes), Validation ($35 \leq t \leq 42$: 34,144 nodes), Test ($43 \leq t \leq 49$: 26,074 nodes).
    \item Automated unit tests in \texttt{tests/test\_no\_leakage.py} confirm zero temporal data leakage.
    \item The complete graph is assembled as a PyTorch Geometric (\texttt{PyG}) \texttt{Data} object and persisted to \texttt{data/processed/elliptic\_graph.pt}.
\end{itemize}

\subsection{Learnable Fourier Time Encoding (O2)}

The \texttt{LearnableFourierTimeEncoder} class was implemented in \texttt{src/models.py} with the following confirmed properties:
\begin{itemize}
    \item Learnable frequency parameters $\boldsymbol{\omega} \in \mathbb{R}^{64}$ initialised from $\mathcal{N}(0, 0.01)$, declared as \texttt{nn.Parameter} for gradient tracking.
    \item Learnable phase parameters $\boldsymbol{\phi} \in \mathbb{R}^{64}$ initialised to zero.
    \item Encoding output $\Phi(\Delta t) \in \mathbb{R}^{128}$ for any input time delta $\Delta t \geq 0$.
    \item Verified gradient flow: $\frac{\partial \mathcal{L}}{\partial \boldsymbol{\omega}} \neq 0$ for all non-zero $\Delta t$ in gradient checks.
    \item Integrated into \texttt{TGATConv} as the temporal component of query, key, and value projections.
    \item \texttt{TemporalNeighborSampler} implements causal masking: only neighbours with $t_j \leq t_i$ are sampled, with up to $K = 20$ neighbours ordered by recency.
\end{itemize}

\subsection{Multi-Pattern TGAT Architecture (O3 — Preliminary)}

The \texttt{MultiPatternTGAT} model has been implemented with the following architecture:
\begin{itemize}
    \item Two-layer TGATConv encoder: $172 \to 128 \to 128$ with 4 attention heads and dropout.
    \item Three parallel classification heads (Circular, Layering, Smurfing), each consisting of a two-layer FFN with ReLU activation and sigmoid output.
    \item Joint multi-task weighted binary cross-entropy loss (Equation~\ref{eq:joint_loss}).
    \item Pattern ground-truth labels mined using Tarjan's SCC (Circular), temporal directed BFS with $\text{length} \geq 3$ and $t_{k+1} \geq t_k$ (Layering), and degree thresholding $d_{\text{out}} \geq 10 \lor d_{\text{in}} \geq 10$ (Smurfing).
\end{itemize}

\noindent \textbf{Preliminary Training Observations:}
\begin{itemize}
    \item The Layering head achieves a validation AUC-ROC of approximately 0.58 in early training epochs, with precision improving as positive-class weighting is tuned.
    \item The Smurfing head achieves a validation AUC-ROC of approximately 0.96 in early training epochs, reflecting the strong discriminative signal provided by the degree-based features (out-degree, fan-out ratio).
    \item The Circular Transfer head produces near-zero recall, consistent with the Bitcoin DAG domain constraint documented in Section~\ref{subsubsec:pattern_heads}.
\end{itemize}

\noindent \textbf{Note:} Full training convergence results, hyperparameter-tuned performance, and extended baseline comparison metrics are deferred to Dissertation Phase~2.

Figure~\ref{fig:auprc_comparison} presents the comparative evaluation of precision-recall dynamics (AUPRC) across graph architectures on the chronological test split.

\begin{figure}[!htbp]
\centering
\includegraphics[width=0.88\textwidth]{auprc_comparison.png}
\caption[AUPRC Comparison Across Graph Architectures]{Area Under the Precision-Recall Curve (AUPRC) comparison across baseline graph models and TemporalAML on the Elliptic chronological test split ($t \in [43, 49]$). Given the severe 2.2\% illicit class imbalance, AUPRC provides a critical performance measure for practical AML triage by penalising false-positive alert volume.}
\label{fig:auprc_comparison}
\end{figure}

\subsection{GNNExplainer Integration (O4 — Implemented)}

The \texttt{src/explain.py} module implements the GNNExplainer pipeline with the following confirmed components:
\begin{itemize}
    \item Per-node, per-pattern explanation generation: for each flagged node $v^*$ with $\hat{y}^{p}_{v^*} > 0.5$, the pattern-specific head's prediction is explained.
    \item Edge mask optimisation using 200 gradient ascent steps (Adam optimiser, learning rate 0.01).
    \item Top-$k$ edge selection ($k = 10$ by default) from the optimised mask $\mathbf{M}^*$.
    \item Time-ordered evidence subgraph construction: selected nodes sorted by timestamp $t$.
    \item Automated SAR narrative template generation (JSON format).
    \item Integration with the Streamlit forensic dashboard (\texttt{app/dashboard.py}): interactive subgraph visualisation with edge time delta labels and mask confidence scores.
\end{itemize}

---

\section{Discussion}
\label{sec:discussion}

\subsection{Importance of Temporal Modelling}

The Phase~1 implementation confirms that incorporating continuous temporal information into graph attention fundamentally changes the information that the model can use. The learnable Fourier encoder projects transaction time differences $\Delta t$ into a 128-dimensional space, allowing the attention coefficient $\alpha_{ij}$ (Equation~\ref{eq:attention}) to be sensitive to both the structural similarity of node features \textit{and} the temporal distance between transactions. This is a qualitative improvement over static GAT, which would assign equal temporal importance to a transaction hop occurring one minute ago versus one occurring 6 months ago.

The strong early AUC-ROC achieved by the Smurfing head (approximately 0.96) relative to the Layering head (approximately 0.58) in preliminary training suggests that degree-based structural features (engineered in O1) are highly discriminative for fan-out detection, while layering detection benefits more from temporal velocity signals that require more training epochs to calibrate through the Fourier frequency parameters.

\subsection{Importance of Learnable Frequency Encoding}

The gradient derivation in Equation~\ref{eq:fourier_gradient} demonstrates that frequency parameters $\boldsymbol{\omega}$ receive non-zero gradients for any non-zero time delta. This theoretical property ensures that the Fourier encoding module can adapt its temporal resolution to the specific transaction velocity distributions present in the Elliptic dataset. Preliminary gradient monitoring during training confirms that $\|\frac{\partial \mathcal{L}}{\partial \boldsymbol{\omega}}\|$ is non-negligible and evolving across training steps—direct evidence that the frequencies are being calibrated by the laundering-relevant temporal patterns in the training data.

\subsection{Multi-Pattern Detection Design Rationale}

The shared-embedding multi-task architecture (one TGAT backbone, three heads) is preferred over three independent models for the following reasons:
\begin{itemize}
    \item \textbf{Parameter efficiency:} A single 2-layer TGAT encoder serves all three detection tasks, reducing the total parameter count by approximately 67\% compared to training three separate models.
    \item \textbf{Cross-pattern feature sharing:} Laundering patterns are not mutually exclusive. A transaction participating in a layering chain may simultaneously exhibit smurfing characteristics. Shared representations allow the model to learn cross-pattern co-occurrence signals.
    \item \textbf{Training stability:} Multi-task learning with positive-class-weighted BCE provides gradient regularisation, reducing overfitting on the rare illicit class.
\end{itemize}

\subsection{Explainability for AML Compliance}

Attention weights alone—while providing some interpretability—are insufficient as AML evidence for two reasons: (i) they are not guaranteed to be faithful to the model's actual decision process (attention is not explanation~\cite{jain2019attention}); and (ii) they do not directly produce a subgraph that compliance officers can trace. GNNExplainer's mutual information optimisation (Equation~\ref{eq:gnnexplainer}) provides a principled, post-hoc mechanism that identifies the minimal causal subgraph sufficient to reproduce the model's classification—a directly actionable output for SAR preparation.

\subsection{Current Phase 1 Limitations}

\begin{itemize}
    \item \textbf{Circular Transfer detection:} Due to the Bitcoin UTXO DAG constraint, no transaction-level cycles exist in the Elliptic graph. This architectural inclusion is retained for applicability to wallet-entity-level graphs but is expected to yield negligible results on the current dataset.
    \item \textbf{Incomplete training convergence:} The full multi-task model has not yet converged to optimal performance in Phase~1. Hyperparameter tuning, learning rate scheduling, and focal loss integration are planned for Phase~2.
    \item \textbf{Preliminary metrics:} The AUC-ROC and loss curves reported above are from early training runs and should not be interpreted as final performance metrics.
    \item \textbf{Explainability quality not yet quantified:} Fidelity, sparsity, and SAR usability metrics for the GNNExplainer outputs are pending Phase~2 evaluation.
\end{itemize}

\noindent \textbf{All full quantitative performance metrics (F1-score, AUC-ROC, Precision, Recall, AUPRC, ablation results, and baseline comparisons) are to be evaluated in Dissertation Phase~2.}
